\section{Background and Related Work}\label{sec:related}

\paragraph{Ontology portability.}
Gruber~\cite{gruber1993} framed portability of knowledge bases as translation: systems commit to
a shared ontology and translators map it into each system's representation. The ontology fixes
meaning, and the translators stay small because they deal only in syntax and conventions. Our
adapters play the role of those translators, with one change: the systems on either side are
LLMs. The write-in direction (text to ontology) is therefore an extraction problem rather than a
syntactic mapping, and it is where we measure most of the loss. Industrial interoperability
frameworks~\cite{ontointerop2025} and LLM-based ontology matching such as
Agent-OM~\cite{qiang2024agentom} address alignment \emph{between} ontologies. We avoid alignment
by having all engines commit to one small ontology, and we discuss when that is not enough in
Sect.~\ref{sec:discussion}.

\paragraph{Agent communication semantics.}
Surveys of agent interoperability protocols~\cite{ehtesham2025survey} describe MCP, A2A and
related protocols as layers of transport and message syntax. A semantic reading of these
protocols~\cite{beyondmessage2026} separates communication, syntactic and semantic layers and
observes that current protocols standardise the first two. Memory portability sits in the third
layer: two engines can exchange well-formed JSON and still disagree on whether an item is a
current fact about the user or a superseded fact about someone else.

\paragraph{Memory for LLM agents.}
MemGPT~\cite{packer2023memgpt} manages a hierarchy of context and external storage for a single
agent. Mem0~\cite{chhikara2025mem0} extracts and consolidates salient facts from conversations.
Zep~\cite{rasmussen2025zep} stores memory as a temporal knowledge graph with validity intervals,
and AriGraph~\cite{anokhin2024arigraph} combines semantic and episodic memory in a graph for
agents acting in text environments. These systems show that structured, time-aware memory helps
a single agent. None of them treats the memory as an artifact that must keep its meaning when
the engine behind the agent changes, and their schemas are internal to each system. MCP servers
that expose a knowledge graph as tools are close to our read-out layer; what we add is a
published, versioned ontology behind the tools, and a measurement of what survives migration.

\paragraph{Benchmarks.}
LongMemEval~\cite{wu2025longmemeval} evaluates the long-term memory of chat assistants over long
histories of sessions, including knowledge updates and temporal reasoning. We reuse its insight
that updates and abstention must be probed explicitly, but build our own personas: LongMemEval has
no notion of migrating memory between engines, of erasure, or of a structured store, and its
free-text answers lack the competitor strings (superseded values, other people's values) that our
deterministic scoring needs.

\paragraph{Standards and regulation.}
We reuse W3C vocabularies rather than define new ones where they exist: PROV-O~\cite{provo} for
derivation and invalidation, SKOS~\cite{skos} for controlled values, FOAF~\cite{foaf} for people,
DCMI Terms~\cite{dcterms} for metadata and OWL-Time~\cite{owltime} for validity. The W3C AI Agent
Memory Interoperability Community Group~\cite{w3cmemcg2026} was chartered in July 2026; our
ontology and metric are offered as input to such efforts. Article~20 of the GDPR~\cite{gdpr}
gives data subjects a right to receive their data ``in a structured, commonly used and
machine-readable format''; Quebec's Law~25~\cite{quebec25} added a comparable right, in force since September 2024.
Neither says what format would make memory usable by a different engine.
